% ===================================================================
%  QUADRO N.º 15 — O Mercado. — Os Alimentos.
%
%  Ceci n'est PAS une transcription : c'est une traduction, faite en
%  2026, en portugais d'aujourd'hui. Voir texto/pt/10-quadro-01.tex
%  pour la consigne, qui vaut pour les seize fichiers.
% ===================================================================

%%K t15-apar-1 apar
\VUsaut{4.0mm}
\VUtitre{15.0pt}{60}{QUADRO N.º 15}
\VUsaut{3.0mm}

%%K t15-tit-1 sub
\VUpk{11.4pt}{40}{O Mercado. --- Os Alimentos.}
\VUsaut{3.0mm}

%%K t15-01-1 p
1. --- A maior parte dos comerciantes que nos fornecem os alimentos
necessários estão reunidos sob a \VUgras{nave} \textsuperscript{(1)} do
\VUgras{mercado coberto}, ou nas ruas vizinhas.

%%K t15-02-1 p
2. --- O \VUgras{salsicheiro} \textsuperscript{(2)}, que vende \VUgras{presunto}
\textsuperscript{(3)}, \VUgras{morcela} \textsuperscript{(4)}, \VUgras{salpicão}
\textsuperscript{(5)}, \VUgras{linguiça} \textsuperscript{(6)} e \VUgras{toucinho}
\textsuperscript{(7)}, está ao lado da \VUgras{vendedeira de manteiga e queijo}
\textsuperscript{(8)}, cuja banca está coberta de \VUgras{queijos de bola}
\textsuperscript{(9)} de casca vermelha, \VUgras{camemberts} \textsuperscript{(10)},
grandes \VUgras{rodas de gruyère} \textsuperscript{(11)} e \VUgras{pães de
manteiga} \textsuperscript{(12)} fresca.

%%K t15-03-1 p
3. --- Diante dela, uma \VUgras{hortaliceira} \textsuperscript{(13)} oferece aos
seus clientes belos \VUgras{tomates} \textsuperscript{(14)}, \VUgras{couves-flores}
\textsuperscript{(15)}, \VUgras{alfaces} \textsuperscript{(16)}, \VUgras{couves}
\textsuperscript{(17)}, \VUgras{abóboras} \textsuperscript{(19)}, \VUgras{melões}
\textsuperscript{(18)} e até \VUgras{cogumelos} \textsuperscript{(20)}. Mas um
\VUgras{cervejeiro}, que parou justamente diante da sua banca o seu
\VUgras{carro de entregas} \textsuperscript{(21)} carregado de \VUgras{barris de
cerveja} \textsuperscript{(22)}, impede que os seus clientes se aproximem. Um
hortelão traz-lhe um cesto de \VUgras{alcachofras} \textsuperscript{(23)}.

%%K t15-tit-2 sub
\VUsaut{4.0mm}
\VUpk{10.2pt}{40}{Vender e comprar.}
\VUsaut{3.0mm}

%%K t15-04-1 p
4. --- No mercado a animação é grande, pois chegou a hora da
\VUgras{arrematação} \textsuperscript{(24)}. As \VUgras{donas de casa}
\textsuperscript{(26)} que vêm aqui fazer as compras detêm-se de passagem diante
da banca da \VUgras{vendedeira de miudezas} \textsuperscript{(25)} para comprar
\VUgras{tripas} ou uma \VUgras{cabeça de vitela}.

%%K t15-05-1 p
5. --- Um \VUgras{cozinheiro} \textsuperscript{(27)} gordo regateia um belíssimo
\VUgras{atum} com a \VUgras{peixeira} \textsuperscript{(28)}, que sobe os preços à
sua vontade. Recebeu uma grande partida de \VUgras{enguias, linguados,
trutas} e \VUgras{carpas}. O seu \VUgras{bacalhau} e os seus
\VUgras{mexilhões} estão muito frescos.

%%K t15-tit-3 sub
\VUsaut{4.0mm}
\VUpk{10.2pt}{40}{O barato e o caro.}
\VUsaut{3.0mm}

%%K t15-06-1 p
6. --- A \VUgras{vendedeira de aves} \textsuperscript{(29)}, instalada ao seu lado,
é muito honrada. Quando vende um \VUgras{peru}, um \VUgras{ganso}, um
\VUgras{pato} ou um simples \VUgras{frango}, não pede mais do que o preço
justo.

%%K t15-07-1 p
7. --- Na esquina da praça, no primeiro andar, uma \VUgras{vendedeira de
roupa branca} \textsuperscript{(30)} estabeleceu-se há pouco, onde os compradores
estão sempre certos de encontrar um sortido muito bom de \VUgras{toalhas de
mesa}, \VUgras{guardanapos, aventais} e \VUgras{panos de cozinha}. No mesmo
andar, um \VUgras{cuteleiro} \textsuperscript{(31)} montou a sua oficina. Afia as
\VUgras{facas} das hortaliceiras e faz aos hortelãos \VUgras{podadeiras} do
\VUgras{aço} mais duro.

%%K t15-tit-4 sub
\VUsaut{4.0mm}
\VUpk{10.2pt}{40}{A mercearia.}
\VUsaut{3.0mm}

%%K t15-08-1 p
8. --- Por baixo da sua oficina abriu-se não há muito uma grande loja de
comestíveis. Já não é a \VUgras{mercearia} \textsuperscript{(32)} simples e modesta
de outros tempos, que só vendia géneros correntes --- \VUgras{chá}
\textsuperscript{(33)}, \VUgras{chocolate} \textsuperscript{(34)}, \VUgras{pão de açúcar}
\textsuperscript{(37)} e \VUgras{sabão} \textsuperscript{(38)}. Aqui vende-se de tudo:
\VUgras{bolachas estaladiças}, \VUgras{conservas} \textsuperscript{(35)},
\VUgras{licores} \textsuperscript{(36)}, \VUgras{rum, conhaque, kirsch, anis} e
\VUgras{curaçau}. Encontra-se aqui caça com a mesma facilidade ---
\VUgras{lebre} \textsuperscript{(39)}, \VUgras{tordos} \textsuperscript{(40)},
\VUgras{perdizes} \textsuperscript{(41)}, \VUgras{codornizes} \textsuperscript{(42)},
\VUgras{patos bravos} \textsuperscript{(43)} ou \VUgras{faisão} \textsuperscript{(44)} ---
como fruta tropical, tal como \VUgras{bananas} \textsuperscript{(46)} e
\VUgras{ananases}.

%%K t15-09-1 p
9. --- Um \VUgras{caixeiro} \textsuperscript{(45)} muito prestativo está disposto a
entregar o que se lhe peça: \VUgras{doce} \textsuperscript{(47)}, de groselha ou
de marmelo, \VUgras{banha} \textsuperscript{(48)}, \VUgras{pepininhos}
\textsuperscript{(49)} ou \VUgras{sardinhas} \textsuperscript{(50)} de salmoura.

%%K t15-09-2 p
Isso é muito cómodo para as \VUgras{clientes} \textsuperscript{(51)}, que
encontram, ao lado dos géneros mais correntes, as primícias mais raras e a
fruta mais saborosa: \VUgras{peras} \textsuperscript{(52)} sumarentas,
\VUgras{ameixas} \textsuperscript{(53)}, \VUgras{uvas} \textsuperscript{(54)},
\VUgras{pêssegos} \textsuperscript{(55)}, \VUgras{amêndoas} \textsuperscript{(56)} e
\VUgras{nozes} \textsuperscript{(57)}.

%%K t15-tit-5 sub
\VUsaut{4.0mm}
\VUpk{10.2pt}{40}{Os clientes.}
\VUsaut{3.0mm}

%%K t15-10-1 p
10. --- O \VUgras{arroz} \textsuperscript{(58)} e as \VUgras{lentilhas}
\textsuperscript{(59)} estão junto à caixa de \VUgras{arenques} \textsuperscript{(60)}
fumados; a \VUgras{batata} \textsuperscript{(61)} e o \VUgras{feijão}
\textsuperscript{(63)} acotovelam-se com as \VUgras{maçãs} \textsuperscript{(62)}, os
\VUgras{figos} \textsuperscript{(64)}, as \VUgras{ameixas secas} \textsuperscript{(65)} e
as \VUgras{avelãs} \textsuperscript{(66)}. Também se encontram nesta loja
\VUgras{ovos} \textsuperscript{(67)} frescos e \VUgras{azeitonas}
\textsuperscript{(68)}, de modo que os \VUgras{cestos} \textsuperscript{(70)} e as
\VUgras{redes das compras} \textsuperscript{(73)} se enchem a toda a pressa.

%%K t15-11-1 p
11. --- O \VUgras{café} \textsuperscript{(72)} desta excelente casa ganhou uma fama
bem merecida entre as \VUgras{criadas} \textsuperscript{(69)} e as cozinheiras,
pois o velho \VUgras{merceeiro} \textsuperscript{(71)} torra-o ele mesmo com o
maior esmero.

%%K t15-tit-6 sub
\VUsaut{4.0mm}
\VUpk{10.2pt}{40}{Coisas que se veem.}
\VUsaut{3.0mm}

%%K t15-12-1 p
12. --- Nem toda a gente vem ao mercado comprar provisões. No pitoresco ir
e vir da ruela que leva à nave, vê-se a gente mais variada. Junto a um
bonacheirão \VUgras{magarefe} \textsuperscript{(75)} que empurra diante de si um
\VUgras{carrinho de mão} \textsuperscript{(74)}, estão dois soldados de licença,
muito orgulhosos de ostentar os seus vistosos uniformes: o \VUgras{hussardo}
\textsuperscript{(76)} com o seu \VUgras{dólmã} azul e a sua \VUgras{barretina de
penacho}, e o \VUgras{cabo de zuavos} com os seus calções largos e um
\VUgras{fez} na cabeça.

%%K t15-13-1 p
13. --- O \VUgras{padeiro} \textsuperscript{(80)}, de pé à soleira da
\VUgras{padaria} \textsuperscript{(78)}, acaba de amassar a massa do seu
\VUgras{pão} \textsuperscript{(79)} e de tirar do forno os \VUgras{croissants}
\textsuperscript{(81)}, os \VUgras{pãezinhos} \textsuperscript{(82)} e as \VUgras{broas}
\textsuperscript{(83)} que estão na montra.

%%K t15-14-1 p
14. --- Por cima da padaria mora uma hábil \VUgras{costureira}
\textsuperscript{(84)} que emprega várias \VUgras{bordadeiras} \textsuperscript{(85)}. Ao
seu lado mora uma \VUgras{lavadeira e engomadeira} \textsuperscript{(86)}, que
engoma a \VUgras{roupa branca} \textsuperscript{(88)} depois de a lavar e secar, e
a passa com um \VUgras{ferro de engomar} \textsuperscript{(87)}.

%%K t15-tit-7 sub
\VUsaut{4.0mm}
\VUpk{10.2pt}{40}{Carne, peixe e legumes.}
\VUsaut{3.0mm}

%%K t15-15-1 p
15. --- No \VUgras{talho} \textsuperscript{(89)} do rés do chão vende-se carne de
todas as espécies --- \VUgras{carneiro} \textsuperscript{(90)}, \VUgras{vaca}
\textsuperscript{(94)} ou \VUgras{vitela} \textsuperscript{(92)} --- em \VUgras{pernas de
cordeiro, bifes, assados} e \VUgras{costeletas}. O \VUgras{ajudante de
talho} \textsuperscript{(93)} corta toda esta carne e põe de parte os \VUgras{ossos
com tutano} \textsuperscript{(96)}. À porta do talho está uma \VUgras{vendedeira de
ostras} \textsuperscript{(97)} que vende \VUgras{ostras} \textsuperscript{(98)}. Abre-as
se lho pedirem.

%%K t15-16-1 p
16. --- Todos estes comerciantes pagam uma patente ou um direito de
banca; por isso o \VUgras{polícia} \textsuperscript{(100)} persegue sem piedade as
pobres \VUgras{hortaliceiras ambulantes} \textsuperscript{(99)}, que lhes fazem
concorrência levando nos seus carrinhos \VUgras{cenouras, rabanetes, nabos,
alhos-porros} ou \VUgras{molhos de espargos, ervilhas frescas, espinafres,
azedas, agriões} e \VUgras{salsa}.

%%K t15-tit-8 sub
\VUsaut{4.0mm}
\VUpk{10.2pt}{40}{Cuidado com o polícia!}
\VUsaut{3.0mm}

%%K t15-17-1 p
17. --- O rapaz que vende \VUgras{alhos} \textsuperscript{(101)} e \VUgras{cebolas}
sabe muito bem que a ele também não lhe é permitido vender a sua
mercadoria perto do mercado. Deita a correr, com o risco de virar o cesto
de \VUgras{caranguejos}, \VUgras{lagostins} e \VUgras{camarões} da
\VUgras{vendedeira} \textsuperscript{(102)} de \VUgras{lagostas} e
\VUgras{lavagantes}.

%%K t15-18-1 p
18. --- Toda a gente se afadiga e faz uma grande algazarra; mas isso não
impede que se ouça claramente o pregão do \VUgras{vidraceiro}
\textsuperscript{(103)} ambulante, nem o apregoar do \VUgras{carvoeiro}
\textsuperscript{(104)}, que acaba de deixar um momento o seu carro para entregar
um \VUgras{saco de carvão} \textsuperscript{(105)}.
\VUsaut{6.0mm}
\VUfilet{22mm}
